\documentclass[12pt]{article}
\usepackage[utf8]{inputenc}
\usepackage[spanish]{babel}
\usepackage{listings}
\usepackage{color}
\usepackage{graphicx}
\usepackage{geometry}
\usepackage{hyperref}
\geometry{margin=2.5cm}

\title{SEGUIMIENTO V}
\author{Sara Rondon, Natalia Ángel\\ Electrónica digital II}
\date{}

\definecolor{codegray}{rgb}{0.95,0.95,0.95}
\definecolor{mygreen}{rgb}{0,0.6,0}

\lstset{
    backgroundcolor=\color{codegray},
    basicstyle=\ttfamily\footnotesize,
    breaklines=true,
    frame=single,
    captionpos=b,
    numbers=left,
    numberstyle=\tiny,
    keywordstyle=\color{blue},
    commentstyle=\color{mygreen},
    stringstyle=\color{red},
    showstringspaces=false
}

\begin{document}
\maketitle
\tableofcontents
\newpage

\section{Identificación del proyecto}
\begin{itemize}
    \item \textbf{Nombre del proyecto:} Sistema IoT de monitoreo de temperatura, humedad y movimiento
    \item \textbf{Integrantes del equipo:} Sara Rondon, Natalia Ángel
    \item \textbf{Asignatura:} Electrónica Digital II
\end{itemize}

\section{Resumen}
El proyecto consiste en el desarrollo de un sistema IoT basado en el microcontrolador ESP32, orientado al monitoreo de variables ambientales y de movimiento en tiempo real. El sistema permite medir temperatura y humedad mediante un sensor DHT, así como detectar el estado dinámico del sistema mediante un acelerómetro MPU6050.

Se implementa un servidor web embebido que permite visualizar los datos en tiempo real desde un navegador, así como un bot de Telegram que permite el monitoreo remoto y la interacción mediante comandos. El sistema también incluye un conjunto de alertas multimodales que combinan señales sonoras, visualización web y notificaciones remotas.

\section{Descripción del hardware}

El sistema está basado en el microcontrolador ESP32, encargado del procesamiento de datos y la comunicación del sistema.

\textbf{Entradas}
\begin{itemize}
    \item Sensor DHT11/DHT22 (GPIO 4): medición de temperatura y humedad.
    \item Sensor MPU6050 (I2C GPIO 21 SDA, GPIO 22 SCL): medición de aceleración y detección de movimiento.
    \item Pulsador (GPIO 18): botón de pánico.
\end{itemize}

\textbf{Procesamiento}
\begin{itemize}
    \item Lectura de sensores y cálculo de variables ambientales.
    \item Análisis de movimiento mediante magnitud del vector de aceleración.
    \item Evaluación de umbrales y generación de alertas.
\end{itemize}

\textbf{Salidas}
\begin{itemize}
    \item Buzzer pasivo (GPIO 27): generación de alertas sonoras.
    \item Comunicación web mediante WiFi.
    \item Notificaciones remotas mediante bot de Telegram.
\end{itemize}

El sistema es alimentado mediante el puerto USB del ESP32.

\begin{figure}
        \centering
        \includegraphics[width=0.5\linewidth]{DIAGRAMA5.png}
        \caption{Diagrama de conexión del sistema}
\end{figure}

\section{Descripción del software}
El software está desarrollado en MicroPython y utiliza programación asíncrona mediante la librería \texttt{uasyncio}.

El sistema se divide en múltiples tareas concurrentes que permiten ejecutar simultáneamente la lectura de sensores, el servidor web y la comunicación con Telegram. Esto garantiza un funcionamiento eficiente y en tiempo real.

Se utilizan librerías como \texttt{network} para conexión WiFi, \texttt{urequests} para comunicación HTTP y \texttt{ujson} para manejo de datos.

\section{Estructura del código}
El código se encuentra organizado en los siguientes bloques:

\begin{itemize}
    \item Configuración de parámetros globales y umbrales.
    \item Inicialización del hardware (sensores, buzzer y pulsador).
    \item Estado global del sistema.
    \item Funciones auxiliares.
    \item Sistema de alertas.
    \item Tareas asíncronas (sensores, servidor web y Telegram).
    \item Programa principal.
\end{itemize}

\section{Explicación de funciones principales}

\begin{itemize}
    \item \texttt{obtener\_ip():} Obtiene la dirección IP del ESP32.
    \item \texttt{url\_encode():} Codifica mensajes para envío mediante URL.
    \item \texttt{enviar\_telegram():} Envía mensajes al bot de Telegram.
    \item \texttt{alerta\_sonora():} Genera diferentes patrones sonoros según el tipo de alerta.
    \item \texttt{disparar\_alerta():} Gestiona las alertas evitando envíos repetitivos.
    \item \texttt{leer\_dht\_seguro():} Lee el sensor DHT manejando errores.
    \item \texttt{tarea\_sensores():} Gestiona la lectura de sensores y evaluación de condiciones.
    \item \texttt{servidor\_web():} Implementa el servidor HTTP embebido.
    \item \texttt{tarea\_telegram():} Gestiona la interacción con el bot de Telegram.
\end{itemize}

\section{Comunicación (si aplica)}
La comunicación del sistema se realiza mediante conexión WiFi.

El ESP32 ejecuta un servidor web embebido que permite la visualización de datos en tiempo real. Adicionalmente, se implementa un bot de Telegram que permite monitorear el sistema de forma remota y recibir alertas.

\section{Interfaz de usuario (si aplica)}
La interfaz de usuario está compuesta por dos elementos principales:

\textbf{Interfaz web}
\begin{itemize}
    \item Visualización de temperatura, humedad y movimiento.
    \item Estado de alarmas.
    \item Actualización automática en tiempo real.
\end{itemize}

\textbf{Bot de Telegram}
\begin{itemize}
    \item Permite consultar variables del sistema.
    \item Permite recibir alertas.
    \item Incluye comandos como /estado, /temp, /humedad y /movimiento.
\end{itemize}

\section{Manejo de errores y seguridad}
El sistema utiliza estructuras \texttt{try/except} para garantizar su estabilidad frente a errores de lectura de sensores o fallos de comunicación.

Además, se implementa control de spam en las alertas y liberación de memoria mediante \texttt{gc.collect()}, asegurando un funcionamiento continuo sin bloqueos.

\section{Código fuente}
\begin{lstlisting}
import gc
import time
import ujson
import dht
import uasyncio as asyncio
import network
import urequests

from machine import Pin, I2C, PWM
from mpu6050 import MPU6050


# ======================================================
# CONFIGURACION GENERAL
# ======================================================

# IMPORTANTE:
TIPO_DHT = "DHT11"

# Configuracion de Telegram.
# Cambiar por los datos reales del bot.
BOT_TOKEN = "8714231259:AAFgi3bfxKQXWop5iThKQC0tMV80P_EPDMM"
CHAT_ID = "8889804535"

# Pines ESP32.
PIN_DHT = 4
PIN_SDA = 21
PIN_SCL = 22
PIN_BUZZER = 27
PIN_PULSADOR = 18

# Umbrales ambientales.
TEMP_MIN = 18.0
TEMP_MAX = 35.0
HUM_MIN = 40.0
HUM_MAX = 70.0

# Umbrales de movimiento en unidades de g.
MOV_NORMAL_G = 1.15
MOV_BRUSCO_G = 1.80

# Intervalos de operacion.
INTERVALO_SENSORES = 2       # segundos
INTERVALO_TELEGRAM = 5       # segundos
ANTI_SPAM_ALERTA = 60        # segundos entre alertas repetidas del mismo tipo


# ======================================================
# INICIALIZACION DE HARDWARE
# ======================================================

if TIPO_DHT.upper() == "DHT11":
    sensor_dht = dht.DHT11(Pin(PIN_DHT))
else:
    sensor_dht = dht.DHT22(Pin(PIN_DHT))

i2c = I2C(0, sda=Pin(PIN_SDA), scl=Pin(PIN_SCL), freq=400000)
mpu = MPU6050(i2c)

buzzer = PWM(Pin(PIN_BUZZER))
buzzer.duty(0)

pulsador = Pin(PIN_PULSADOR, Pin.IN, Pin.PULL_UP)


# ======================================================
# ESTADO GLOBAL
# ======================================================

estado = {
    "temperatura": None,
    "humedad": None,
    "ax": 0.0,
    "ay": 0.0,
    "az": 0.0,
    "magnitud_g": 1.0,
    "movimiento": "SIN DATO",
    "alarma": "INICIANDO",
    "panico": "INACTIVO",
    "ip": "SIN IP",
    "dht_ok": False,
    "ultima_actualizacion": "SIN DATO"
}

ultimo_envio_alerta = {}
telegram_offset = 0


# ======================================================
# FUNCIONES AUXILIARES
# ======================================================

def obtener_ip():
    wlan = network.WLAN(network.STA_IF)
    if wlan.isconnected():
        return wlan.ifconfig()[0]
    return "SIN CONEXION WIFI"


def url_encode(texto):
    """
    Codificacion basica para enviar texto por URL a Telegram.
    Evita caracteres problematicos en MicroPython.
    """
    texto = str(texto)

    reemplazos = {
        "%": "%25",
        " ": "%20",
        "\n": "%0A",
        ":": "%3A",
        "/": "%2F",
        "#": "%23",
        "&": "%26",
        "?": "%3F",
        "=": "%3D",
        "+": "%2B",
        "°": "%C2%B0",
        "¡": "",
        "!": "%21",
        "á": "a", "é": "e", "í": "i", "ó": "o", "ú": "u",
        "Á": "A", "É": "E", "Í": "I", "Ó": "O", "Ú": "U",
        "ñ": "n", "Ñ": "N"
    }

    for original, codificado in reemplazos.items():
        texto = texto.replace(original, codificado)

    return texto


def telegram_configurado():
    return (
        BOT_TOKEN != "TU_TOKEN_DE_TELEGRAM"
        and CHAT_ID != "TU_CHAT_ID"
        and len(BOT_TOKEN) > 10
        and len(CHAT_ID) > 0
    )


def mensaje_estado():
    temp = "SIN LECTURA" if estado["temperatura"] is None else "{:.1f} C".format(estado["temperatura"])
    hum = "SIN LECTURA" if estado["humedad"] is None else "{:.1f} %".format(estado["humedad"])

    return (
        "ESTADO DEL SISTEMA\n"
        "Temperatura: {}\n"
        "Humedad: {}\n"
        "DHT: {}\n"
        "Movimiento: {}\n"
        "Magnitud: {:.2f} g\n"
        "Alarma: {}\n"
        "Boton panico: {}\n"
        "IP: {}"
    ).format(
        temp,
        hum,
        "OK" if estado["dht_ok"] else "SIN LECTURA",
        estado["movimiento"],
        estado["magnitud_g"],
        estado["alarma"],
        estado["panico"],
        estado["ip"]
    )


def mensaje_umbrales():
    return (
        "UMBRALES CONFIGURADOS\n"
        "Temperatura minima: {:.1f} C\n"
        "Temperatura maxima: {:.1f} C\n"
        "Humedad minima: {:.1f} %\n"
        "Humedad maxima: {:.1f} %\n"
        "Movimiento normal: > {:.2f} g\n"
        "Movimiento brusco: > {:.2f} g"
    ).format(TEMP_MIN, TEMP_MAX, HUM_MIN, HUM_MAX, MOV_NORMAL_G, MOV_BRUSCO_G)


async def enviar_telegram(mensaje):
    """
    Envia mensaje por Telegram.
    Si BOT_TOKEN o CHAT_ID no estan configurados, solo imprime en consola.
    """
    if not telegram_configurado():
        print("Telegram no configurado. Mensaje local:", mensaje)
        return False

    try:
        url = (
            "https://api.telegram.org/bot{}/sendMessage?chat_id={}&text={}"
            .format(BOT_TOKEN, CHAT_ID, url_encode(mensaje))
        )

        respuesta = urequests.get(url)
        codigo = respuesta.status_code
        respuesta.close()

        if codigo == 200:
            print("Mensaje enviado a Telegram.")
            return True
        else:
            print("Telegram respondio con codigo:", codigo)
            return False

    except Exception as e:
        print("Error enviando Telegram:", e)
        return False


async def alerta_sonora(tipo):
    """
    Patrones sonoros diferenciados.
    """
    if tipo == "TEMPERATURA":
        frecuencia = 1000
        repeticiones = 2
    elif tipo == "HUMEDAD":
        frecuencia = 1300
        repeticiones = 3
    elif tipo == "MOVIMIENTO":
        frecuencia = 2000
        repeticiones = 4
    elif tipo == "PANICO":
        frecuencia = 2500
        repeticiones = 5
    elif tipo == "DHT":
        frecuencia = 700
        repeticiones = 1
    else:
        frecuencia = 900
        repeticiones = 1

    for _ in range(repeticiones):
        buzzer.freq(frecuencia)
        buzzer.duty(512)
        await asyncio.sleep(0.16)
        buzzer.duty(0)
        await asyncio.sleep(0.12)

    buzzer.duty(0)


async def disparar_alerta(tipo, mensaje):
    """
    Activa alerta sonora y alerta remota evitando spam.
    """
    ahora = time.time()
    ultimo = ultimo_envio_alerta.get(tipo, 0)

    await alerta_sonora(tipo)

    if ahora - ultimo >= ANTI_SPAM_ALERTA:
        ultimo_envio_alerta[tipo] = ahora
        await enviar_telegram("ALERTA {}:\n{}".format(tipo, mensaje))

    print("ALERTA", tipo, mensaje)


def leer_dht_seguro():
    """
    Lee el DHT evitando que ETIMEDOUT detenga todo el sistema.
    Si falla, mantiene los valores anteriores y marca dht_ok = False.
    """
    try:
        sensor_dht.measure()
        temperatura = float(sensor_dht.temperature())
        humedad = float(sensor_dht.humidity())

        estado["temperatura"] = temperatura
        estado["humedad"] = humedad
        estado["dht_ok"] = True

        return temperatura, humedad, True

    except Exception as e:
        print("Error leyendo DHT:", e)
        estado["dht_ok"] = False

        # Mantiene ultimos valores validos.
        return estado["temperatura"], estado["humedad"], False


# ======================================================
# TAREA 1: MONITOREO DE SENSORES
# ======================================================

async def tarea_sensores():
    estado["ip"] = obtener_ip()

    await enviar_telegram(
        "ESP32 iniciado correctamente.\nIP del servidor web: http://{}".format(estado["ip"])
    )

    while True:
        alarmas = []

        try:
            # Lectura DHT protegida.
            temperatura, humedad, dht_ok = leer_dht_seguro()

            # Lectura MPU6050 protegida.
            try:
                ax, ay, az = mpu.leer_aceleracion_g()
                magnitud = mpu.leer_magnitud_g()
                movimiento = mpu.clasificar_movimiento(magnitud)

                estado["ax"] = ax
                estado["ay"] = ay
                estado["az"] = az
                estado["magnitud_g"] = magnitud
                estado["movimiento"] = movimiento

            except Exception as e:
                print("Error leyendo MPU6050:", e)
                estado["movimiento"] = "ERROR MPU6050"
                magnitud = estado["magnitud_g"]
                movimiento = estado["movimiento"]

            estado["panico"] = "ACTIVO" if pulsador.value() == 0 else "INACTIVO"
            estado["ultima_actualizacion"] = str(time.time())
            estado["ip"] = obtener_ip()

            # Solo evalua temperatura y humedad si el DHT tiene lectura valida.
            if dht_ok:
                if temperatura < TEMP_MIN:
                    alarmas.append("TEMPERATURA BAJA")
                    await disparar_alerta(
                        "TEMPERATURA",
                        "Temperatura baja: {:.1f} C. Limite minimo: {:.1f} C"
                        .format(temperatura, TEMP_MIN)
                    )

                if temperatura > TEMP_MAX:
                    alarmas.append("TEMPERATURA ALTA")
                    await disparar_alerta(
                        "TEMPERATURA",
                        "Temperatura alta: {:.1f} C. Limite maximo: {:.1f} C"
                        .format(temperatura, TEMP_MAX)
                    )

                if humedad < HUM_MIN:
                    alarmas.append("HUMEDAD BAJA")
                    await disparar_alerta(
                        "HUMEDAD",
                        "Humedad baja: {:.1f} %. Limite minimo: {:.1f} %"
                        .format(humedad, HUM_MIN)
                    )

                if humedad > HUM_MAX:
                    alarmas.append("HUMEDAD ALTA")
                    await disparar_alerta(
                        "HUMEDAD",
                        "Humedad alta: {:.1f} %. Limite maximo: {:.1f} %"
                        .format(humedad, HUM_MAX)
                    )
            else:
                alarmas.append("DHT SIN LECTURA")
                await disparar_alerta(
                    "DHT",
                    "No se pudo leer el sensor DHT. Revise tipo de sensor, DATA GPIO 4, VCC 3V3 y GND."
                )

            if estado["movimiento"] == "MOVIMIENTO BRUSCO":
                alarmas.append("MOVIMIENTO BRUSCO")
                await disparar_alerta(
                    "MOVIMIENTO",
                    "Movimiento brusco detectado. Magnitud: {:.2f} g"
                    .format(estado["magnitud_g"])
                )

            if pulsador.value() == 0:
                alarmas.append("BOTON DE PANICO")
                await disparar_alerta(
                    "PANICO",
                    "Boton de panico activado manualmente."
                )

            if len(alarmas) == 0:
                estado["alarma"] = "NORMAL"
                buzzer.duty(0)
            else:
                estado["alarma"] = " / ".join(alarmas)

        except Exception as e:
            estado["alarma"] = "ERROR GENERAL"
            buzzer.duty(0)
            print("Error general en tarea_sensores:", e)

        gc.collect()
        await asyncio.sleep(INTERVALO_SENSORES)


# ======================================================
# TAREA 2: SERVIDOR WEB
# ======================================================

def pagina_web():
    color_alarma = "#0f9d58" if estado["alarma"] == "NORMAL" else "#d93025"

    temp_web = "SIN LECTURA" if estado["temperatura"] is None else "{:.1f}".format(estado["temperatura"])
    hum_web = "SIN LECTURA" if estado["humedad"] is None else "{:.1f}".format(estado["humedad"])

    html = """<!DOCTYPE html>
<html>
<head>
    <meta charset="UTF-8">
    <meta http-equiv="refresh" content="3">
    <title>Seguimiento V - ESP32</title>
    <style>
        body {{
            font-family: Arial, sans-serif;
            background: #f3f6fb;
            margin: 0;
            padding: 24px;
            color: #1f2933;
        }}
        .card {{
            max-width: 760px;
            margin: auto;
            background: white;
            border-radius: 14px;
            padding: 24px;
            box-shadow: 0 4px 18px rgba(0,0,0,0.12);
        }}
        h1 {{
            margin-top: 0;
            color: #0b5394;
        }}
        .grid {{
            display: grid;
            grid-template-columns: repeat(2, 1fr);
            gap: 14px;
        }}
        .box {{
            background: #eef3f8;
            border-radius: 10px;
            padding: 14px;
        }}
        .value {{
            font-size: 28px;
            font-weight: bold;
        }}
        .alarm {{
            background: {color_alarma};
            color: white;
            border-radius: 10px;
            padding: 14px;
            font-size: 20px;
            font-weight: bold;
            margin: 16px 0;
        }}
        .small {{
            font-size: 13px;
            color: #5f6b7a;
        }}
        table {{
            width: 100%;
            border-collapse: collapse;
            margin-top: 12px;
        }}
        td, th {{
            padding: 8px;
            border-bottom: 1px solid #d8dee9;
            text-align: left;
        }}
    </style>
</head>
<body>
    <div class="card">
        <h1>Monitoreo IoT Biomedico - ESP32</h1>
        <p class="small">Servidor web embebido con actualizacion automatica cada 3 segundos.</p>

        <div class="alarm">Estado de alarma: {alarma}</div>

        <div class="grid">
            <div class="box">
                <div>Temperatura</div>
                <div class="value">{temperatura} C</div>
            </div>
            <div class="box">
                <div>Humedad relativa</div>
                <div class="value">{humedad} %</div>
            </div>
            <div class="box">
                <div>Estado de movimiento</div>
                <div class="value">{movimiento}</div>
            </div>
            <div class="box">
                <div>Magnitud aceleracion</div>
                <div class="value">{magnitud:.2f} g</div>
            </div>
        </div>

        <h2>Estado del sistema</h2>
        <table>
            <tr><th>Variable</th><th>Valor</th></tr>
            <tr><td>Sensor DHT</td><td>{dht_estado}</td></tr>
            <tr><td>Boton de panico</td><td>{panico}</td></tr>
            <tr><td>Aceleracion X</td><td>{ax:.2f} g</td></tr>
            <tr><td>Aceleracion Y</td><td>{ay:.2f} g</td></tr>
            <tr><td>Aceleracion Z</td><td>{az:.2f} g</td></tr>
            <tr><td>IP ESP32</td><td>{ip}</td></tr>
        </table>

        <h2>Umbrales configurados</h2>
        <table>
            <tr><th>Parametro</th><th>Rango / limite</th></tr>
            <tr><td>Temperatura</td><td>{tmin:.1f} C a {tmax:.1f} C</td></tr>
            <tr><td>Humedad</td><td>{hmin:.1f} % a {hmax:.1f} %</td></tr>
            <tr><td>Movimiento</td><td>Normal &gt; {mov:.2f} g</td></tr>
            <tr><td>Movimiento brusco</td><td>&gt; {brusco:.2f} g</td></tr>
        </table>
    </div>
</body>
</html>
""".format(
        color_alarma=color_alarma,
        alarma=estado["alarma"],
        temperatura=temp_web,
        humedad=hum_web,
        movimiento=estado["movimiento"],
        magnitud=estado["magnitud_g"],
        dht_estado="OK" if estado["dht_ok"] else "SIN LECTURA",
        panico=estado["panico"],
        ax=estado["ax"],
        ay=estado["ay"],
        az=estado["az"],
        ip=estado["ip"],
        tmin=TEMP_MIN,
        tmax=TEMP_MAX,
        hmin=HUM_MIN,
        hmax=HUM_MAX,
        mov=MOV_NORMAL_G,
        brusco=MOV_BRUSCO_G
    )

    return html


async def servidor_web():
    async def atender_cliente(reader, writer):
        try:
            request = await reader.read(1024)
            request = request.decode()

            if "GET /json" in request:
                contenido = ujson.dumps(estado)
                writer.write("HTTP/1.1 200 OK\r\n")
                writer.write("Content-Type: application/json\r\n")
                writer.write("Connection: close\r\n\r\n")
                writer.write(contenido)
            else:
                contenido = pagina_web()
                writer.write("HTTP/1.1 200 OK\r\n")
                writer.write("Content-Type: text/html\r\n")
                writer.write("Connection: close\r\n\r\n")
                writer.write(contenido)

            await writer.drain()
            await writer.wait_closed()

        except Exception as e:
            print("Error cliente web:", e)

    await asyncio.start_server(atender_cliente, "0.0.0.0", 80)
    print("Servidor web iniciado en http://{}".format(obtener_ip()))

    while True:
        await asyncio.sleep(3600)


# ======================================================
# TAREA 3: BOT DE TELEGRAM
# ======================================================

async def tarea_telegram():
    global telegram_offset

    if not telegram_configurado():
        print("Telegram no configurado. Configure BOT_TOKEN y CHAT_ID para activar el bot.")
        return

    while True:
        try:
            url = (
                "https://api.telegram.org/bot{}/getUpdates?offset={}"
                .format(BOT_TOKEN, telegram_offset)
            )

            respuesta = urequests.get(url)
            datos = respuesta.json()
            respuesta.close()

            if datos.get("ok"):
                for item in datos.get("result", []):
                    telegram_offset = item["update_id"] + 1

                    mensaje = item.get("message", {})
                    texto = mensaje.get("text", "").strip().lower()

                    if texto == "/estado":
                        await enviar_telegram(mensaje_estado())

                    elif texto == "/temp":
                        if estado["temperatura"] is None:
                            await enviar_telegram("Temperatura: SIN LECTURA")
                        else:
                            await enviar_telegram(
                                "Temperatura actual: {:.1f} C".format(estado["temperatura"])
                            )

                    elif texto == "/humedad":
                        if estado["humedad"] is None:
                            await enviar_telegram("Humedad: SIN LECTURA")
                        else:
                            await enviar_telegram(
                                "Humedad actual: {:.1f} %".format(estado["humedad"])
                            )

                    elif texto == "/movimiento":
                        await enviar_telegram(
                            "Movimiento: {}\nMagnitud: {:.2f} g"
                            .format(estado["movimiento"], estado["magnitud_g"])
                        )

                    elif texto == "/umbrales":
                        await enviar_telegram(mensaje_umbrales())

                    elif texto == "/ayuda":
                        await enviar_telegram(
                            "Comandos disponibles:\n"
                            "/estado\n"
                            "/temp\n"
                            "/humedad\n"
                            "/movimiento\n"
                            "/umbrales\n"
                            "/ayuda"
                        )

                    elif texto:
                        await enviar_telegram("Comando no reconocido. Use /ayuda")

        except Exception as e:
            print("Error revisando Telegram:", e)

        gc.collect()
        await asyncio.sleep(INTERVALO_TELEGRAM)


# ======================================================
# PROGRAMA PRINCIPAL
# ======================================================

async def main():
    print("Iniciando sistema Seguimiento V...")
    estado["ip"] = obtener_ip()

    asyncio.create_task(tarea_sensores())
    asyncio.create_task(servidor_web())
    asyncio.create_task(tarea_telegram())

    while True:
        await asyncio.sleep(10)


try:
    asyncio.run(main())
finally:
    buzzer.duty(0)
    asyncio.new_event_loop()
\end{lstlisting}

\end{document}

